\section{Introduction}
\label{sec:intro}

Deep neural networks are vulnerable to adversarial examples---inputs crafted 
with small, imperceptible perturbations that cause confident misprediction 
\cite{szegedy2014intriguing,goodfellow2015explaining}.
Despite more than a decade of research, the arms race between attacks and 
defenses remains unresolved: architecturally-coupled defenses require model 
retraining, are often architecture-specific, and may degrade accuracy on clean 
data \cite{croce2020reliable,athalye2018obfuscated}.

A complementary direction is \emph{runtime detection}: rather than hardening 
the model, intercept adversarial inputs before they affect predictions.
Existing detectors often provide no formal statistical guarantees on false 
alarm rates, rely on features that are themselves vulnerable to adaptive 
attacks, or require knowledge of the model's training procedure 
\cite{ma2018characterizing,grosse2017statistical,lee2018simple}.

\textbf{Key insight.} 
Clean and adversarial inputs traverse fundamentally different paths through 
a neural network's activation manifold.
While the output may look confident in both cases, the \emph{topology} of 
intermediate representations---captured by persistent homology---differs 
systematically \cite{gebhart2019adversary,carlini2019evaluating}.
\prism exploits this topological signal at runtime as a 
computationally tractable anomaly score.

\textbf{Contributions.}
We make four contributions:

\begin{enumerate}[leftmargin=*, label=\textbf{C\arabic*}]

\item \textbf{Topological activation profiler (\tamm).}
We compute per-layer persistence features relative to a clean medoid profile
and combine them with DCT and softmax-entropy features in a calibrated
ensemble. The ensemble-no-TDA ablation measures the marginal value of the
topological channel.

\item \textbf{Conformal detection guarantee (\cadg).}
We apply split conformal prediction to the anomaly score stream, obtaining 
\emph{distribution-free} FPR certificates at three severity tiers 
(L1: $\alpha{=}0.10$; L2: $\alpha{=}0.03$; L3: $\alpha{=}0.005$)
using only a held-out clean calibration set. Conformal prediction has been
applied to outlier and out-of-distribution detection before---iDECODe
\cite{kaur2022idecode} calibrates an equivariance nonconformity score and
Bates et al.\ \cite{bates2023outliers} give conformal $p$-value outlier
tests---but for \emph{novelty/OOD} inputs in input or logit space. To our
knowledge \prism is the first to attach such a distribution-free
false-positive-rate certificate to an \emph{adversarial-input} detector
built on multi-layer activation topology. The guarantee is exact under
exchangeability; we additionally \emph{audit} it under benign covariate
shift (CIFAR-10-C) and report the measured FPR inflation rather than a
qualitative caveat (\S\ref{sec:cert_shift}).

\item \textbf{Temporal campaign detector (\sacd).}
We monitor the score stream with a CUSUM-with-drift-floor change detector
and report \emph{time-to-detect} (mean queries between the first adversarial
step and the first L0 trigger) as the headline metric. On sustained
$\rho{=}1.0$ streams \sacd raises the campaign alert in $2.6 \pm 0.5$
queries while maintaining $0.0$ false-alarm activation on clean-only
streams---a threat model absent from prior work. Under a surrogate-free
direct-score attack that drives the single-shot tier down, \sacd flags the
probing campaign and holds the full-system miss to $0.03$
(\S\ref{sec:direct_score}), so campaign detection is load-bearing for
adaptive robustness, not an auxiliary feature.

\item \textbf{Routed MoE recovery (\tamsh).}
For inputs reaching the L3 severity tier, \prism routes the input through
an attack-specialist expert sub-network rather than issuing a blanket
rejection. A \emph{learned} router---logistic regression on the pooled
final-layer activation, trained on a held-out L3 pool and evaluated on a
disjoint split (5 seeds)---recovers $28.7\%$ of L3-rejected PGD
adversarials, $9.1\times$ the uniform $1/K$ control ($3.2\%$) and
$2.1\times$ the lightweight $H_0{+}H_1$ Wasserstein gate ($13.7\%$),
matching a prior-augmented upper bound ($29.3\%$) that hardcodes the PGD
specialist---without that hardcoding. The large margin over the uniform
control isolates the contribution to the routing signal rather than
expert capacity
(Appendix~\ref{app:tamsh_diag}).

\end{enumerate}

Empirically, we report locked multi-seed ($5{\times}n{=}1000$) artifact
results across three full CNN settings --- CIFAR-10/ResNet-18,
CIFAR-10/WRN-28-10, and CIFAR-100/ResNet-18 --- covering detection
(FGSM, PGD, Square, CW, and AutoAttack where run), four reproduced
baselines (LID, Mahalanobis, ODIN, Energy), campaign streams, recovery,
and a per-channel ablation that exposes each feature family's
attack-specific failure mode. We additionally report a standard-attack
CIFAR-10/ViT-B/16 transfer check for FGSM, PGD-50, and Square only. On the
main CIFAR-10/ResNet-18 setting, \prism Full outperforms the strongest
baseline (Mahalanobis) by $+33.6$~pp mean TPR at matched FPR;
WRN-28-10 holds a $+16.5$~pp gap over Mahalanobis and CIFAR-100 retains
$+42.6$~pp over ODIN (Table~\ref{tab:baselines_multi}). The full
CIFAR-100 per-component breakdown is in Appendix~\ref{app:cifar100}.
The ensemble-complete adaptive PGD claim is scoped to CIFAR-10/ResNet-18;
WRN and CIFAR-100 adaptive probes are appendix stress tests with narrower
coverage.

\textbf{Scope and limitations.}
\prism is a \emph{runtime monitor}, not an adversarial trainer; 
it complements (but does not replace) robust training.
Detection latency scales with the number of monitored layers and the TDA
subsample size; the submitted results report measured GPU latency from the
same artifact pipeline used for detection.
Like all detectors, \prism does not provide a guarantee against 
adaptive adversaries who are aware of the topological scoring mechanism.
